\section{Feedback variables and frequency two}\label{sec:frequency}
\subsection{A universal repair law}
The next argument works on the original factor scopes, including on boxes
with positive lower endpoints.

\begin{lemma}[Feedback repair and gluing]\label{lem:feedback-repair}
Fix singleton Bernoulli means and a set $F$ of $f_F$ variable vertices
whose deletion leaves the factor incidence graph acyclic. For any nonnegative local functions $g_e$ and any prescribed-mean
local laws $P_e$ on their scopes, there is one global prescribed-mean law
$P'$ satisfying
\[
 \E_{P'}g_e\ge2^{-f_F}\E_{P_e}g_e\qquad\hbox{for every }e.
\]
\end{lemma}
\begin{proof}
Extend each local law to $F\cup e$ by independently sampling any missing
feedback coordinates. Put $C=2^{f_F}$. Construct a global reference law
$Q$ with a common uniform $U$: each feedback coordinate independently
chooses, fairly, its success interval $[0,x_j]$ or $[1-x_j,1]$;
every outside coordinate uses $[0,x_i]$. Write $Q_F$ for its feedback
marginal and $Q_i$ for its $(F,i)$ marginal. For a feedback state $s$
and outside state $b$, every prescribed-mean local law satisfies
\[
 P(s,b)\le m(s,b;i):=\min\bigl(\{x_j:s_j=1\},
 \{1-x_j:s_j=0\},\{x_i\text{ if }b=1;\ 1-x_i\text{ if }b=0\}\bigr).
\]
The minimum includes all feedback coordinates and the outside coordinate.
Choose the feedback orientations that place every desired event at the
left endpoint if $b=1$, and at the right endpoint if $b=0$. This single
orientation pattern has probability $1/C$, and its intersection has
length $m(s,b;i)$. Therefore $Q_i(s,b)\ge P(s,b)/C$.
The same endpoint argument gives $Q_F(s)\ge P_F(s)/C$, also when there
are no outside variables. For $F=\varnothing$ its unique state has mass
one.

For a fixed extended factor law $P_e$, define the nonnegative residuals
\[
 h_s=Q_F(s)-(P_e)_F(s)/C,\qquad
 r_{i,s,b}=Q_i(s,b)-(P_e)_{F,i}(s,b)/C\quad(i\in e\setminus F).
\]
For each $i$, $\sum_b r_{i,s,b}=h_s$. At state $s$ with $h_s>0$, place
residual mass $h_s$ and sample these outside coordinates independently
with probabilities $r_{i,s,b}/h_s$. For $h_s=0$ all these residuals
vanish. For an empty outside scope place just the mass $h_s$.
The resulting measure $R_e$ has total mass $1-1/C$, and
$P'_e=P_e/C+R_e$ has feedback marginal $Q_F$ and $(F,i)$ marginal $Q_i$,
while dominating $P_e/C$ entrywise.

Condition on any $s$ with $Q_F(s)>0$. The factor laws $P'_e(\cdot\mid s)$
have consistent outside singleton marginals $Q_i(s,b)/Q_F(s)$. Their
incidence graph is a forest. Root each nontrivial component at a factor,
sample its law, then traverse the tree. Each newly reached factor shares
exactly one already sampled variable; sample its new variables using its
local conditional law given that variable. Null conditioning states are
never encountered. Isolated variables use their stated conditional means;
factors with no outside variables need no sampling. This preserves every
conditional factor law. Mix over $s$ with weights $Q_F(s)$. The resulting
global law has marginals $P'_e$ on $F\cup e$, and nonnegativity gives
the claimed domination of expectations. There is one loss $C$, regardless
of the lengths of the forest paths.
\end{proof}

\begin{theorem}[Feedback-variable gap]\label{thm:feedback}
For positive monomials on any finite nonnegative box, deletion of $f_F$
variable vertices to an incidence forest implies $\tgap\le2^{f_F}\hgap$.
\end{theorem}
\begin{proof}
Substitute fixed coordinates first. In normalized endpoint variables $Y$,
an original physical product is $\varphi_e(Y)=\sum_{A\subseteq e}
c_A\prod_{i\in A}Y_i$ with $c_A\ge0$. This expansion is used only
inside the original scope. For each nonempty $A$ choose an index $k_A$
of smallest normalized mean, and let
\[
 \ell_e(Y)=c_\varnothing+\sum_{\varnothing\ne A\subseteq e}c_AY_{k_A}.
\]
On binary vertices $\ell_e\ge\varphi_e$, while every prescribed-mean
law has $\E\ell_e=\cav\varphi_e$, by common upper attainment.
Thus $g_e=\ell_e-\varphi_e$ is a nonnegative function on the original
scope whose local maximum expectation is the original factor's exact
gap. Apply Lemma~\ref{lem:feedback-repair} to maximizing local laws,
then sum with the positive coefficients. The full upper value is the
sum of local upper values, proving the result. Replacing the factors by
expanded hyperedges would change the incidence graph and would not
justify this proof.
\end{proof}

For $f_F=0$ this recovers classical incidence-forest exactness; exactness
of standard linearization on Berge-acyclic hypergraphs is established
in \citet{DelPiaKhajavirad2018}. The general domination lemma is sharp
for arbitrary nonnegative local payoffs: on $f_F$ feedback bits and one
outside bit, all fair, use the $2^{f_F+1}$ indicators of their states.
Every local maximum is $1/2$, while their sum is identically one.
Private variables fixed to one can make the scopes distinct. These
indicators have zero literals, so they do not prove sharpness for positive
monomials when $f_F>1$.

\begin{proposition}[One feedback variable is sharp]\label{prop:feedback-sharp}
For $n\ge2$, let $f(a,x)=a\sum_{i=1}^nx_i+\prod_{i=1}^nx_i$ and evaluate
at $a=1/n$, $x_i=1-1/n$ on the unit cube. Then
$\tgap=2-1/n$ and $\hgap=1$. Its incidence graph has treewidth two and
becomes a tree after deleting $a$.
\end{proposition}
\begin{proof}
The bilinear terms contribute one in total and the large product
contributes $(n-1)/n$. If $R$ counts failed leaves and $A$ is the
binary anchor, the full deficiency is
$\E[AR]+\Prob(R\ge1)-1/n$, with $\E R=1$ and $\E A=1/n$.
Pointwise $AR+\ind\{R\ge1\}\le R+A$, giving the upper bound one.
Exactly one uniformly chosen failed leaf and an independent anchor
attain equality. Write $e_0$ for the large factor and $e_i$ for $ax_i$.
Path bags $\{a,e_0,x_i\}$ with attached bags $\{a,x_i,e_i\}$ give width
two. A cycle rules out width one, and deleting $a$ leaves a tree.
\end{proof}

\subsection{Half-integral degree slabs}
For frequency at most two, form a loopless dual multigraph with factor
vertices. A shared variable is an edge between its two factors; a private
variable is an edge to its own new dummy vertex. Parallel edges are
allowed. Unused variables may be completed independently. Dummy factors
have zero weight or identically zero costs. Let $g$ be its odd girth,
with $g=\infty$ in the bipartite case.

\begin{lemma}[Degree-slab vertices]\label{lem:degree-slab}
For a loopless multigraph and $p\in[0,1]^E$, put
$s_v(p)=\sum_{i\in\delta(v)}p_i$. Every vertex of
\[
 \mathcal S(p)=\{z\in[0,1]^E:
       \lfloor s_v(p)\rfloor\le s_v(z)\le\lceil s_v(p)\rceil\ (v\in V)\}
\]
has fractional edges forming vertex-disjoint odd cycles, all valued $1/2$.
\end{lemma}
\begin{proof}
Let $J$ be the fractional edges at a vertex and $R$ the tight degree rows
incident to them, counting a repeated lower/upper equality once.
The restricted rows must have full column rank on $J$: otherwise a
small perturbation in a null direction preserves every tight constraint
and box bound, contradicting extremality. Hence $|J|\le|R|$.
An integer tight residual degree cannot contain exactly one fractional
edge, so every such row has at least two fractional incidences. Thus
$2|R|\le2|J|$. Equality forces each fractional edge to have both endpoints
in $R$, and each such vertex to have fractional degree two. These are
cycles; an even cycle admits an alternating perturbation and is excluded.
On an odd cycle the two incident fractional values sum to the integer
one, forcing all values to $1/2$. The argument excludes parallel
fractional two-cycles and degree-one dummy vertices automatically.
\end{proof}
This is the usual fractional matching rank argument, here stated for
integer degree slabs; compare the fixed-degree formulation in
\citet[Theorem 2.1]{Barrus2014}.

\begin{theorem}[Frequency two and odd girth]\label{thm:frequency-two}
For positive monomials of frequency at most two on a unit cube,
\[
 \tgap\le\frac{g}{g-1}\hgap\quad(g<\infty),
 \qquad \tgap=\hgap\quad(g=\infty).
\]
Every finite odd-girth constant is sharp, and the unrestricted
frequency-two supremum is $3/2$.
\end{theorem}
\begin{proof}
Use failure probabilities $p_i=1-x_i$, and define
$b_v(p)=\max_{i\in\delta(v)}p_i$ and $c_v(p)=\min\{1,s_v(p)\}$,
with empty values zero. For a failure edge set $J$, write $I_v(J)$ for
its coverage indicator. Exact single-product envelopes give
\begin{equation}\label{eq:coverage-baseline}
 \tgap=\sum_v a_v(c_v-b_v),\qquad
 \hgap=\max_{\Prob(i\in J)=p_i}\sum_v a_v(\E I_v(J)-b_v).
\end{equation}
In particular an approximation of unshifted coverage alone is insufficient.
Decompose $p=\E Z$ into vertices of $\mathcal S(p)$. On each allowed
slab $c_v$ is affine (linear below one, constant above), so
$\E c_v(Z)=c_v(p)$; convexity gives $\E b_v(Z)\ge b_v(p)$.

Given $Z=z$, keep integral edges. On each fractional odd cycle of length
$l$, choose one of its $l$ maximum matchings uniformly, then fairly
choose that matching or its complement in the cycle. Each edge has
marginal $1/2$. A matching leaves its uniformly selected vertex uncovered;
its complement covers every vertex. Thus coverage probability is
$1-1/(2l)$ at each cycle vertex. A vertex with another selected integral
edge is covered surely. Otherwise its $c_v(z)=1$ and $b_v(z)=1/2$.
All vertices off these cycles have equal integral $b_v(z)=c_v(z)$.
With $\alpha=1-1/g$, this proves
\[
 \E[I_v(J)\mid Z]\ge\alpha c_v(Z)+(1-\alpha)b_v(Z).
\]
Averaging retains the baseline:
$\E I_v(J)-b_v(p)\ge\alpha[c_v(p)-b_v(p)]$.
Sum this in \eqref{eq:coverage-baseline}. For a bipartite graph there
are no fractional cycles, so the slab decomposition itself attains
$c_v(p)$ simultaneously, proving exactness.

For sharpness take a unit-weight odd cycle of length $l$ with all edge
means $1/2$. Its termwise gap and coverage baseline both equal $l/2$.
For every selected edge set $J$,
$\sum_vI_v(J)\le|J|+(l-1)/2$: if $|J|\le(l-1)/2$ use $2|J|$,
otherwise use $l$. Expected coverage is therefore at most $l-1/2$,
attained by the matching/complement law. The hull gap is $(l-1)/2$.
\end{proof}
Bipartite equality is a statement about the scalar envelopes. It does
not assert compatibility of all separately feasible product-value
coordinates in the full lifted multilinear polytope. Zero-lower boxes
transfer by scaling; unequal positive aspect ratios need a separate analysis.

\subsection{Convex cardinality factors}
A discrete sequence $\varphi(0),\ldots,\varphi(d)$ is convex if
$\varphi(k+2)-2\varphi(k+1)+\varphi(k)\ge0$. Its cardinality factor
is the unique multiaffine interpolant of the values
$\varphi(\sum_iX_i)$ on binary vertices. It is generally not the literal
continuous composition $\varphi(\sum_ix_i)$. Sequences may be negative,
decreasing, or nonmonotone; only their discrete convexity is needed.

\begin{lemma}[Cardinality envelopes]\label{lem:cardinality-envelopes}
Let $f_e$ be such a factor and $\widehat\varphi_e$ its piecewise-affine
interpolation between integers. Then
\[
 \vex f_e(p)=\widehat\varphi_e\Bigl(\sum_{i\in e}p_i\Bigr).
\]
A prescribed-mean law supported on the two adjacent cardinalities attains
this minimum. One common-threshold law maximizes every such factor and
their sum. If $p_{(1)}\ge\cdots\ge p_{(d)}$ are the local means, its
upper value is
\[
 \varphi_e(0)+\sum_{j=1}^d[\varphi_e(j)-\varphi_e(j-1)]p_{(j)}.
\]
\end{lemma}
\begin{proof}
Jensen's inequality gives the lower bound for the convex function
$\widehat\varphi_e$. The cube slab between the two adjacent integer
sums is integral: two fractional coordinates admit an opposite
perturbation; one fractional coordinate cannot make an integer sum
bound tight and admits a small perturbation. Decompose $p$ into its
binary vertices to attain equality.
Discrete convexity implies supermodularity of the cardinality set
function. Among laws maximizing a factor, choose one maximizing also
$\E|X|^2$. If two incomparable support sets have positive mass, replace
equal mass on them by their intersection and union. Singleton means
are unchanged; supermodularity cannot decrease the primary value,
while the secondary value strictly increases by twice the product of
the two nonempty set differences' sizes. This is impossible. The
maximizer is chain-supported, and the chain law with prescribed means
is precisely the common-threshold law. The same law therefore maximizes
each factor individually, and summing establishes simultaneous upper
attainment. Integrating its cardinality gives the displayed formula.
\end{proof}
The single-factor principles are classical; symmetric envelope formulas
appear in \citet{Sherali1997}, and the common-aspect product formula in
\citet[Proposition 4.1]{AdamsGupteXu2019}.

\begin{theorem}[Cardinality odd-girth bound]\label{thm:cardinality-frequency}
For a sum of convex cardinality factors of frequency at most two,
with each entire factor kept in the relaxation, the sharp constants of
Theorem~\ref{thm:frequency-two} hold unchanged.
\end{theorem}
\begin{proof}
Use success probabilities in the degree slab, and decompose $p=\E Z$
as before. Each local lower envelope is affine on its allowed slab, so
\[
 \E\vex f_v(Z)=\vex f_v(p),\qquad \E\tgap(Z)\le\tgap(p),
\]
where the second assertion also uses concavity of each upper envelope.
At a fractional-cycle vertex $v$, let $m$ count its incident integral
ones. Its local lower value is $\varphi_v(m+1)$, and its upper value
is $[\varphi_v(m)+\varphi_v(m+2)]/2$, since precisely two coordinates
are half-valued. Matching/complement rounding on a length-$l$ cycle
selects zero or two of them with probability $1/(2l)$ each, and one
otherwise. Therefore the expected local cost is exactly
$\vex f_v(Z)+\tgap_v(Z)/l$. Off the cycles the local gap is zero.
Averaging gives
\[
 \vex f(p)\le\E f(X)\le\sum_v\vex f_v(p)+\tgap(p)/g.
\]
Subtract from the common upper value. Bipartite slabs are integral and
need no rounding. Positive bilinear odd cycles prove sharpness.
\end{proof}

\begin{corollary}[Common-aspect physical products]\label{cor:cardinality-box}
For original positive products on boxes $[\ell_i,\rho_e\ell_i]$,
$\ell_i>0$, assume each scope $e$ has a common coordinate aspect ratio
$\rho_e>1$. Then the frequency-two odd-girth theorem applies to the
original factors. Its constant is sharp for every fixed common ratio.
\end{corollary}
\begin{proof}
A normalized binary vertex gives the table
$a_e(\prod_{i\in e}\ell_i)\rho_e^K$, which is discrete convex.
Shared variables enforce ratio compatibility, so different connected
components may have different ratios. No new factors are introduced.
For a bilinear odd cycle on $[1,\rho]^n$, expansion gives affine terms
plus $(\rho-1)^2$ times the unit-cube bilinear polynomial. Both gaps
scale by the same positive number. Fixed coordinates are substituted
first, and ratio one is trivial.
\end{proof}

\paragraph{Formal verification scope.}
The \href{../formal/topics/19-structural-multilinear/COVERAGE.md}{topic 19 proofs}
cover the feedback and frequency-two gap statements above from their original
incidence assumptions, including the cardinality and stated box extensions.
The feedback-one flower has minimum feedback size one and treewidth exactly
two; the odd-cycle witnesses have the required frequency and exact odd girth.
The generic-payoff sharpness proof also includes distinct scopes with private
mean-one variables. It does not establish positive-monomial sharpness for
arbitrary feedback size. The unequal-aspect ratio $7/6$ is verified, while a
universal frequency-two bound of $3/2$ on arbitrary positive boxes remains
open here. The
\href{../formal/topics/19-structural-multilinear/VERIFICATION.md}{verification record}
identifies the checks and source snapshots. The algorithms and bit-complexity
claims in the following subsection are outside this Lean scope.

\subsection{Rational scalar-envelope algorithms}
\begin{proposition}[Frequency-two binary oracle]\label{prop:edge-cover-oracle}
For a rational positive polynomial of frequency at most two on a unit
cube and arbitrary rational $\lambda$, one can minimize
$f(X)-\lambda^{\mathsf T}X$ over binary $X$ in polynomial bit time.
\end{proposition}
\begin{proof}
With $Z=1-X$, the objective is
\[
 -\sum_i\lambda_i+\sum_i\lambda_iZ_i+
                     \sum_v a_v\ind\{v\text{ uncovered by }Z\}.
\]
Dummy vertices have zero penalties. Include every negative-cost edge:
adding it strictly lowers cost and cannot increase an uncovered penalty.
Add its cost to the constant, remove it from the remaining choices, and
set its endpoints' penalties to zero, retaining those vertices. Remaining
costs and penalties are nonnegative. Add a hub $h$, its mate $k$, their
zero-cost edge $hk$, and an edge $vh$ of cost equal to each vertex $v$'s
penalty. Any selected original edge set extends to an edge cover by
paying exactly its uncovered penalties and adding $hk$. Conversely,
retaining original edges of an augmented cover leaves uncovered only
vertices whose penalty edges were paid. These two comparisons prove
value equality with minimum-weight edge cover.
For completeness this reduces to matching: for a nonnegative-cost graph
without isolated vertices let $\mu(v)$ be its cheapest incident cost.
From a matching $M$, add a cheapest edge at each unmatched vertex. Its
cover cost is at most
$\sum_v\mu(v)-\sum_{uv\in M}(\mu(u)+\mu(v)-c_{uv})$.
Conversely an inclusion-minimal cover is a union of stars (a path of
three edges or a cycle would have a removable edge). Choose one edge
from each star as a matching. Every other edge can be charged to its
unmatched leaf, with $\mu(v)$ no larger than its cost. This shows the
minimum of the displayed expression is no larger than the minimum cover
cost. Hence equality holds and maximum-weight matching supplies an
optimal cover. Parallel edges can keep their cheapest representative.
Weighted matching is polynomial-time \citep{Edmonds1965}.
\end{proof}

\begin{proposition}[Convex-degree matching oracle]\label{prop:cardinality-matching}
For explicitly listed rational discrete-convex tables and arbitrary
rational edge costs $c_i$, the binary problem
\[
 \min_{J\subseteq E}\sum_v\varphi_v(\deg_J(v))+\sum_{i\in J}c_i
\]
is solvable in polynomial bit time.
\end{proposition}
\begin{proof}
At vertex $v$ create $d_v$ slots with nondecreasing charges
$\Delta_v(j)=\varphi_v(j)-\varphi_v(j-1)$. For each edge $i=uv$, create
mandatory vertices $i_u,i_v$, joined by an inactive zero-cost edge.
Connect $i_u$ to all slots of $u$ with cost $\Delta_u(j)+c_i$, and
$i_v$ to all slots of $v$ with cost $\Delta_v(j)$. A matching covering
all mandatory vertices either uses the inactive edge or uses two slots;
mixed activation cannot cover both endpoints. Thus it represents one
binary choice per original edge. Every chosen set is representable,
and the occupied slots at a vertex can be reassigned to its cheapest
$k$ slots because their available gadget neighbors are identical.
Their charges telescope to $\varphi_v(k)-\varphi_v(0)$.

To enforce mandatory coverage by ordinary unconstrained matching, let
$W$ be the sum of absolute original gadget edge costs and subtract
$M=2W+1$ per mandatory endpoint covered by each edge. The all-inactive
matching covers every mandatory vertex. Missing even one such vertex
loses bonus $M$, exceeding any possible original cost difference.
Thus an optimum covers them all, also with negative increments or
negative $c_i$. Adding $\sum_v\varphi_v(0)$ and undoing the fixed bonus
recovers the desired value. The gadget has $O(|E|)$ vertices and
$O(\sum_vd_v^2)$ edges with polynomial rational encoding length.
This is the incremental-slot matching mechanism of convex degree
optimization; compare \citet[Theorem 1.2 and Section 3]{DezaOnn2021}.
\end{proof}

\begin{corollary}[Exact scalar envelopes and supporting minorants]
\label{cor:scalar-envelope-algorithm}
Both preceding factor classes admit exact rational convex-envelope
evaluation, a supporting affine minorant, and scalar graph-hull
separation in polynomial bit time. This includes rational common-aspect
physical products. No polynomial-size explicit description of the full
polytope carrying all individual factor values is asserted.
\end{corollary}
\begin{proof}
The finite vertex-law LP has dual
\[
 \max_{\eta,\lambda}\{\eta+\lambda^{\mathsf T}p:
                 \eta+\lambda^{\mathsf T}z\le f(z)\ (z\in\{0,1\}^n)\}.
\]
Its separation oracle is precisely binary minimization of
$f(z)-\lambda^{\mathsf T}z$, with arbitrary signed linear costs.
Affine summands are absorbed into these costs and restored in the
value; unused coordinates optimize independently. Let
$A=\sum_v\max_j|\varphi_v(j)|$, or $A=\sum_va_v$ for positive monomials,
adding the absolute coefficient sum of separate affine terms if present.
Then $|f(z)|\le A$. The dual polyhedron contains no line, since the rows
$(1,z)$ span $\R^{n+1}$, and its attained optimal face contains a vertex.
Such a vertex solves $n+1$ independent tight equations with a zero-one
matrix. Cramer's rule bounds each coordinate by $(n+1)!A$; its nonzero
integer denominator determinant has magnitude at least one. Clearing
input denominators also gives polynomial bit length for each coordinate.
Thus the box of radius $\max\{1,(n+1)!A\}$ retains an optimum and has
polynomial encoding size. Rational separation-to-optimization
\citep[Theorem~6.4.9]{GLS1988} gives the exact optimum and minorant, also at boundary
means. Every vertex inequality extends to the whole cube by multiaffine
interpolation. Sorting gives the common-threshold upper value and a
supporting upper affine function; together with cube bounds these
separate the scalar graph hull. Rational common-ratio products generate
their convex tables by polynomially many rational multiplications.
\end{proof}
